\metatitle{Schubert calculus}
\metadescription{A computation-focused introduction to Schubert calculus on Grassmannians.}

\section[schubertCalculus]{Schubert calculus}

\defin{Schubert calculus}\label{schubertCalculusDefinition} is the use of
\hyperref[schubertVarietyDefinition]{Schubert varieties} and their intersection
products to solve enumerative problems.  In the most classical case, this means
computing in the cohomology ring of a \hyperref[grassmannian]{Grassmannian}.
The combinatorics is governed by partitions, Young diagrams,
\hyperref[schurS]{Schur polynomials},
\hyperref[schubertCalculusPieri]{Pieri rules},
\hyperref[schubertCalculusGiambelli]{Giambelli determinants}, and
\hyperref[schurLittlewoodRichardson]{Littlewood--Richardson coefficients}.

Good introductions include \name{William Fulton}'s book on Young tableaux
\cite{Fulton1997} and \name{Maria Gillespie}'s tutorial
\cite{Gillespie2019}, which is adapted from her Schubert calculus mini-course
\cite{Gillespie2017SchubertMiniCourse}.
For the flag-variety and polynomial side, see also
\hyperref[schubert]{Schubert polynomials}.  For a broader map of the
geometric objects, see \hyperref[varieties]{varieties in algebraic
combinatorics}.


\section[grassmannian]{Grassmannians}

Let $V$ be an $n$-dimensional complex vector space.  The
\defin{Grassmannian}\label{grassmannianDefinition} $\operatorname{Gr}(k,V)$,
often written
$\operatorname{Gr}(k,n)$, is the space of $k$-dimensional linear subspaces
$W\subseteq V$.

The projective interpretation is often the most concrete.  A $k$-dimensional
linear subspace of $V=\setC^n$ gives a projective $(k-1)$-plane in
$\mathbb{P}^{n-1}$.  Thus $\operatorname{Gr}(2,4)$ parametrizes lines in
$\mathbb{P}^3$, and $\operatorname{Gr}(3,6)$ parametrizes projective planes in
$\mathbb{P}^5$.

\begin{example*}[Lines in projective space]
A point of $\operatorname{Gr}(2,4)$ is a $2$-dimensional vector subspace
$W\subseteq\setC^4$.  After projectivizing, this is a line
$\mathbb{P}(W)\subseteq\mathbb{P}^3$.  Schubert calculus on
$\operatorname{Gr}(2,4)$ therefore answers questions about lines in
projective three-space.
\end{example*}

The dimension is
\[
  \dim \operatorname{Gr}(k,n)=k(n-k).
\]
One way to see this is that a generic $k$-plane can be represented by a
$k\times n$ matrix in row echelon form.  After choosing pivot columns, the
remaining free entries form a $k(n-k)$-dimensional affine chart.


\section[grassmannianSchubertClasses]{Schubert classes on the Grassmannian}

Fix a \hyperref[flagVarieties]{complete flag}
\[
  0=F_0\subset F_1\subset F_2\subset\dotsb\subset F_n=V,
  \qquad \dim F_i=i.
\]
\defin{Schubert classes}\label{grassmannianSchubertClassDefinition}
in $\operatorname{Gr}(k,n)$ are indexed by partitions
$\lambda$ fitting in the $k\times(n-k)$ rectangle:
\[
  n-k\geq \lambda_1\geq \lambda_2\geq\dotsb\geq\lambda_k\geq 0.
\]
The corresponding \defin{Schubert variety}\label{schubertVarietyDefinition} is
\[
  \Omega_\lambda(F_\bullet)
  \coloneqq
  \left\{
    W\in\operatorname{Gr}(k,n):
    \dim(W\cap F_{n-k+i-\lambda_i})\geq i
    \text{ for } 1\leq i\leq k
  \right\}.
\]
Its cohomology class is denoted $\sigma_\lambda$, and its codimension is
$|\lambda|$.

The Schubert classes form an additive basis:
\[
  H^*(\operatorname{Gr}(k,n),\setZ)
  =
  \bigoplus_{\lambda\subseteq (n-k)^k}\setZ\,\sigma_\lambda.
\]
The top class is $\sigma_{(n-k)^k}$, and integration over the Grassmannian
extracts the coefficient of this class.


\subsection[grassmannianSchubertDictionary]{Dictionary for $\operatorname{Gr}(2,4)$}

The Grassmannian $\operatorname{Gr}(2,4)$ has dimension $4$ and parametrizes
lines in $\mathbb{P}^3$.  The partitions fit in a $2\times 2$ box:
\[
  \emptyset,\quad (1),\quad (2),\quad (1,1),\quad (2,1),\quad (2,2).
\]
With respect to a fixed flag in $\setC^4$, the corresponding geometric
conditions are:
\[
\begin{array}{c|l}
\lambda & \text{condition on the line } L\subseteq\mathbb{P}^3\\
\hline
\emptyset & \text{no condition}\\
(1) & L \text{ meets a fixed line}\\
(2) & L \text{ contains a fixed point}\\
(1,1) & L \text{ lies in a fixed plane}\\
(2,1) & L \text{ contains a fixed point and lies in a fixed plane}\\
(2,2) & L \text{ is a fixed line.}
\end{array}
\]
The codimension is the number of boxes in $\lambda$.


\section[schubertCalculusComputationRules]{Computation rules}

For ordinary Grassmannians, multiplication of Schubert classes is computed by
the same Littlewood--Richardson coefficients as multiplication of
\hyperref[schurS]{Schur polynomials}, with the restriction that partitions must
stay inside the $k\times(n-k)$ rectangle:
\[
  \sigma_\lambda\sigma_\mu
  =
  \sum_{\nu\subseteq(n-k)^k}
  c_{\lambda,\mu}^{\nu}\sigma_\nu.
\]
Equivalently, multiply $\schurS_\lambda\schurS_\mu$ and discard or straighten
terms which do not define classes in the Grassmannian.  More precisely, one
uses the usual quotient presentation of $H^*(\operatorname{Gr}(k,n))$, where
the Schur classes outside the rectangle are reduced by the Grassmannian
relations.


\subsection[schubertCalculusPieri]{Pieri rule}

The special class $\sigma_r$ corresponds to a single row of length $r$.  The
Pieri rule says
\[
  \sigma_r\sigma_\lambda
  =
  \sum_\mu \sigma_\mu,
\]
where the sum is over partitions $\mu$ such that $\mu/\lambda$ is a horizontal
strip of size $r$ and $\mu$ still fits in the $k\times(n-k)$ rectangle.
There is a dual version with $\sigma_{1^r}$ and vertical strips.

\begin{example*}[A Pieri computation in $\operatorname{Gr}(3,7)$]
The rectangle is $3\times4$.  Starting from $\lambda=(2,1)$ and multiplying
by $\sigma_2$, we add two boxes with no two in the same column.  The possible
results are
\[
  (4,1),\qquad (3,2),\qquad (3,1,1),\qquad (2,2,1).
\]
Therefore
\[
  \sigma_2\sigma_{2,1}
  =
  \sigma_{4,1}+\sigma_{3,2}+\sigma_{3,1,1}+\sigma_{2,2,1}
  \qquad\text{in } H^*(\operatorname{Gr}(3,7)).
\]
\end{example*}


\subsection[schubertCalculusGiambelli]{Giambelli formula}

The Giambelli formula expresses every Schubert class in terms of the special
classes $\sigma_r$:
\[
  \sigma_\lambda
  =
  \det\left(\sigma_{\lambda_i+j-i}\right)_{1\leq i,j\leq k},
\]
where $\sigma_0=1$ and $\sigma_r=0$ if $r<0$ or $r>n-k$.

\begin{example*}[A Giambelli computation]
In any Grassmannian whose rectangle contains $(2,1)$, we have
\[
  \sigma_{2,1}
  =
  \begin{vmatrix}
    \sigma_2 & \sigma_3\\
    1 & \sigma_1
  \end{vmatrix}
  =
  \sigma_2\sigma_1-\sigma_3.
\]
In $\operatorname{Gr}(2,4)$, the class $\sigma_3$ is zero because the rectangle
has width $2$, so the same formula gives
\[
  \sigma_{2,1}=\sigma_2\sigma_1.
\]
\end{example*}

\begin{example*}[Recovering $\sigma_{1,1}$]
For $\lambda=(1,1)$,
\[
  \sigma_{1,1}
  =
  \begin{vmatrix}
    \sigma_1 & \sigma_2\\
    1 & \sigma_1
  \end{vmatrix}
  =
  \sigma_1^2-\sigma_2.
\]
Thus $\sigma_1^2=\sigma_2+\sigma_{1,1}$.
\end{example*}


\subsection[schubertCalculusLRTruncation]{Littlewood--Richardson truncation}

Littlewood--Richardson multiplication can be used directly, provided we keep
only the classes which fit in the Grassmannian rectangle.

\begin{example*}[A product in $\operatorname{Gr}(2,5)$]
In symmetric functions,
\[
  \schurS_1\schurS_{2,1}
  =
  \schurS_{3,1}+\schurS_{2,2}+\schurS_{2,1,1}.
\]
The rectangle for $\operatorname{Gr}(2,5)$ is $2\times3$, so the term
$\schurS_{2,1,1}$ has too many rows.  Hence
\[
  \sigma_1\sigma_{2,1}
  =
  \sigma_{3,1}+\sigma_{2,2}
  \qquad\text{in } H^*(\operatorname{Gr}(2,5)).
\]
\end{example*}


\section[schubertCalculusGr24]{The ring of $\operatorname{Gr}(2,4)$}

Here is the complete multiplication table for Schubert classes in
$H^*(\operatorname{Gr}(2,4))$.  Products not visible by symmetry are obtained
by commutativity.
\[
\begin{array}{c|cccccc}
\cdot
& 1 & \sigma_1 & \sigma_2 & \sigma_{1,1} & \sigma_{2,1} & \sigma_{2,2}\\
\hline
1
& 1 & \sigma_1 & \sigma_2 & \sigma_{1,1} & \sigma_{2,1} & \sigma_{2,2}\\
\sigma_1
& \sigma_1 & \sigma_2+\sigma_{1,1} & \sigma_{2,1} & \sigma_{2,1} & \sigma_{2,2} & 0\\
\sigma_2
& \sigma_2 & \sigma_{2,1} & \sigma_{2,2} & \sigma_{2,2} & 0 & 0\\
\sigma_{1,1}
& \sigma_{1,1} & \sigma_{2,1} & \sigma_{2,2} & \sigma_{2,2} & 0 & 0\\
\sigma_{2,1}
& \sigma_{2,1} & \sigma_{2,2} & 0 & 0 & 0 & 0\\
\sigma_{2,2}
& \sigma_{2,2} & 0 & 0 & 0 & 0 & 0
\end{array}
\]

For example,
\[
  \sigma_1^2=\sigma_2+\sigma_{1,1},
  \qquad
  \sigma_1^3=2\sigma_{2,1},
  \qquad
  \sigma_1^4=2\sigma_{2,2}.
\]
Since $\sigma_{2,2}$ is the class of a point, integration gives
\[
  \int_{\operatorname{Gr}(2,4)}\sigma_1^4=2.
\]

\begin{example*}[The classical line-count]
The class $\sigma_1$ is the condition that a line in $\mathbb{P}^3$ meets a
fixed line.  Thus $\sigma_1^4=2\sigma_{2,2}$ says:

\[
  \text{There are }2\text{ lines in }\mathbb{P}^3
  \text{ meeting four general fixed lines.}
\]
\end{example*}


\section[schubertCalculusPoincareDuality]{Poincaré duality}

Let $\lambda\subseteq(n-k)^k$.  The Poincaré dual partition is
\[
  \lambda^\vee
  =
  (n-k-\lambda_k,\dotsc,n-k-\lambda_2,n-k-\lambda_1).
\]
Then
\[
  \int_{\operatorname{Gr}(k,n)}
  \sigma_\lambda\sigma_\mu
  =
  \begin{cases}
  1, & \mu=\lambda^\vee,\\
  0, & \text{otherwise.}
  \end{cases}
\]

\begin{example*}
In $\operatorname{Gr}(2,4)$, the dual of $(1)$ is $(2,1)$, so
\[
  \sigma_1\sigma_{2,1}=\sigma_{2,2}.
\]
Geometrically, a general line-meeting condition and a general
point-in-plane condition cut out one line.
\end{example*}


\section[schubertCalculusFlagVariety]{Flag varieties and Schubert polynomials}

Schubert calculus on the \hyperref[completeFlagDefinition]{complete flag
variety} is indexed by permutations
rather than by partitions.  The polynomial representatives are the
\hyperref[schubert]{Schubert polynomials} $\schubert_w(\xvec)$ introduced by
\name{Alain Lascoux} and \name{M.-P. Schützenberger}
\cite{LascouxSchutzenberger1982}.

For the full \hyperref[completeFlagDefinition]{flag variety} in $\setC^3$, with permutations written in one-line
notation, the six Schubert polynomials are
\[
\begin{array}{c|c}
w & \schubert_w(x_1,x_2)\\
\hline
123 & 1\\
213 & x_1\\
132 & x_1+x_2\\
231 & x_1x_2\\
312 & x_1^2\\
321 & x_1^2x_2.
\end{array}
\]
Therefore
\[
  \schubert_{213}\schubert_{132}
  =
  x_1(x_1+x_2)
  =
  \schubert_{312}+\schubert_{231}.
\]
This is the flag-variety analogue of a Schubert-class multiplication
computation.

More generally, the problem of finding a positive combinatorial rule for the
structure constants
\[
  \schubert_u(\xvec)\schubert_v(\xvec)
  =
  \sum_w c_{u,v}^w\schubert_w(\xvec)
\]
is a central open problem; see
\hyperref[schubertLittlewoodRichardson]{Schubert Littlewood--Richardson
coefficients}.
